\section{Measurement Details and Optional Follow-up}

\subsection{Teacher-conditioned masks at a common joint checkpoint}
\label{app:records}

The diagnostic record contains the joint student checkpoint, optimizer
state, teacher identity, routed domain, prompt IDs, sampled response
and mask, loss reduction, and gradient accumulation denominator.
Each teacher-conditioned gradient and proposed step is computed from
the same restored student state. They are diagnostic update supports,
not a unique historical attribution of the cumulative joint delta.

Record raw and clipped gradients separately. Retain the actual optimizer
step and its momentum and decay settings. A zero-gradient optimizer
branch can identify motion caused by stored momentum or weight decay
when those components dominate a support comparison. Prefer FP32
master-weight differences for cumulative updates when available;
converting already rounded BF16 checkpoints to FP32 does not recover
unrecorded updates. Record exclusions, thresholds, relative scales,
tie rules, and per-layer support counts.

Fixed-density overlap comparisons accompany absolute-threshold masks.
A stratified random baseline preserves layer and source-weight-magnitude
bins when feasible. Degenerate or tie-dominated gradients are marked;
ranking a zero vector does not define a meaningful learned subnetwork.
Common-prefix teacher distances are aggregated using the same domain
weights, and distance estimates computed from truncated vocabulary
outputs are labeled as approximations with their retained mass.

\subsection{Optional analysis if the loss comparison shows a stable difference}
\label{app:optional_loss}

These controls are conditional follow-up, not part of the required
PG/top-$k$ experiment. On a fixed student-prefix bank, independently
resample next actions under the current student and compare action
sample counts, an action-independent detached baseline, and an exact
full-vocabulary gradient on a small bank. These known estimator controls
can determine whether an observed parameter effect tracks sampling
variance. A retained-head plus sampled-tail estimator can distinguish
truncation from the full reverse-KL reference
\citep{oh2026vopd,zhang2026emapg,yu2026denseupdates}.

The independent action draws do not extend the original response.
This keeps prefix selection and weights fixed for the local estimator
comparison. Using the original rollout token with its eventual
response-length weight need not satisfy the same conditional sampling
identity. Forward KL, teacher versus student top-$k$, or within-set
renormalization are objective changes and are labeled separately.
These controls are pursued only if they answer a specific difference
observed in the primary comparison.

\subsection{An optional amplitude calibration for teacher distance}
\label{app:teacher_controls}

A simple artificial target can help distinguish update amplitude from
parameter selection. Freeze $p_0$ and $q$ on the same prefixes and set
\begin{equation}
q_\alpha(v\mid h)=
\frac{p_0(v\mid h)^{1-\alpha}q(v\mid h)^\alpha}{Z_\alpha(h)}.
\label{eq:tilt_teacher}
\end{equation}
At the reference student, the full reverse-KL gradient obeys
\begin{equation}
\left.\nabla_\theta D_{\mathrm{KL}}(p_\theta\|q_\alpha)
\right|_{\theta_0}
=\alpha\left.\nabla_\theta D_{\mathrm{KL}}(p_\theta\|q)
\right|_{\theta_0}.
\label{eq:tilt_gradient}
\end{equation}
The log-ratio is an $\alpha$-scaled version of the original log-ratio
plus a prefix-constant normalizer. The constant vanishes under the
score identity. Positive scaling therefore preserves fixed-density
raw-gradient support while
possibly changing absolute-threshold sparsity. This optional
calibration does not replace the real-teacher comparisons and does
not assert that artificial targets have matched downstream capability.

\subsection{Reproducibility records}

Record exact student, teacher, tokenizer, data, and evaluator revisions;
teacher training lineages; prompt schedules and seeds; decoding,
EOS handling, cap, invalid responses, and loss masks; learning-rate
schedule, optimizer moments, precision, clipping, and decay; candidate
sets, probability normalization, teacher/student retained mass, and
advantage handling. Keep per-prompt outputs, source measurements for
every plot, and actual time and memory costs.

Distinguish training seeds, diagnostic batch replicates, and evaluation
resamples. Pairs sharing a teacher are dependent, and prompt-level
uncertainty does not establish population-level teacher relationships.
Historical synthetic GPAS calculations remain analytical artifacts and
are not empirical evidence for these studies.
